The seventh phase of the Coupled Model Intercomparison Project (CMIP7) is underway.
A number of CMIP7 experiments explore future climate change.
Here we provide the greenhouse gas (GHG) boundary conditions, i.e. forcings,
to be used for the experiments which explore future climate change under a range of scenarios until 2100
and their temporal extensions to 2500 (the ScenarioMIP experiments).
The need to produce these GHG forcings on a tight delivery deadline means that an Earth System Model cannot be used to produce them.
Instead, we use a simple climate model (namely the Model for the Assessment of Greenhouse Gas Induced Climate Change, MAGICC) as the basis for these forcings.
As a result, they will inevitably differ from emissions-driven ESM outputs for the same experiment.
The GHG concentration forcings cover a range of potential futures.
% ZNTODO: values and value checks for all of these
CO_2 ranges from X~ppm in year Y in experiment S to X'~ppm in year Y' in experiment S',
CH_4 ranges from X~ppm in year Y in experiment S to X'~ppm in year Y' in experiment S',
and N_2O ranges from X~ppm in year Y in experiment S to X'~ppm in year Y' in experiment S',
while CFC12 is based on WMO and related projections so is assumed to be the same across all experiments.
Compared to CMIP6,
% ZNTODO: values and value checks for all of these
the range for CO_2 is narrower (X to X'~ppm),
the range for CH_4 is TBD (X to X'~ppm),
the range for N_2O is TBD (X to X'~ppm)
and CFC12 projections are [higher/lower] by X (maybe percent) on average.
The changes are driven by changes in the underlying scenarios
and by producing the underlying GHG concentration projections with an emulator (specifically the MAGICC model)
that is in line with AR6 science and includes an estimate of the impact of climate change on CH_4 emissions from wetlands.
The underlying method can and will be used for any future-focussed experiments,
which we illustrate by also presenting the GHG concentrations to be used for PolMIP's VL-CF experiment.
The GHG concentrations are publicly available and in use by CMIP7 modelling teams.
